\documentclass{umk-matematika}

\begin{document}

\maketitlepage
    {Практическая работа №19}
    {Призма. Площадь поверхности и объём}
    {}

\section{Фундамент (1--4 балла)}

\textbf{Задача 1 (1 балл).} В прямой треугольной призме стороны основания равны $3$ см, $4$ см и $5$ см. Высота призмы $10$ см. Найдите площадь боковой поверхности.

\textbf{Задача 2 (2 балла).} В прямой четырёхугольной призме основание — квадрат со стороной $6$ см. Высота призмы $8$ см. Найдите объём призмы.

\textbf{Задача 3 (3 балла).} Правильная четырёхугольная призма имеет сторону основания $5$ см и высоту $10$ см. Найдите площадь полной поверхности.

\textbf{Задача 4 (4 балла).} Объём прямой призмы равен $240$ см³, площадь основания $30$ см². Найдите высоту призмы.

\section{Стены (5--7 баллов)}

\textbf{Задача 5 (5 баллов).} В прямой треугольной призме основание — прямоугольный треугольник с катетами $6$ см и $8$ см. Высота призмы $12$ см. Найдите объём призмы.

\textbf{Задача 6 (6 баллов).} Диагональ правильной четырёхугольной призмы равна $\sqrt{41}$ см, а диагональ боковой грани $5$ см. Найдите объём призмы.

\textbf{Задача 7 (7 баллов).} В правильной шестиугольной призме сторона основания $2$ см, высота $5$ см. Найдите площадь боковой поверхности.

\section{Кровля (8--10 баллов)}

\textbf{Задача 8 (8 баллов).} Бетонная опора моста имеет форму прямой призмы высотой $4$ м. Основание — прямоугольный треугольник с катетами $1{,}5$ м и $2$ м. Найдите объём бетона, необходимый для $3$ таких опор.

\textbf{Задача 9 (9 баллов).} Складское помещение имеет форму прямой призмы с основанием — равнобедренной трапецией. Основания трапеции $8$ м и $4$ м, высота трапеции $3$ м. Длина склада $10$ м. Найдите объём помещения.

\textbf{Задача 10 (10 баллов).} Колонна состоит из правильной четырёхугольной призмы (сторона основания $4$ дм, высота $10$ дм) и куба со стороной $4$ дм, поставленного сверху. Найдите общий объём колонны.

\newpage
\section*{Ответы}

\begin{enumerate}
  \item $120$ см²
  \item $288$ см³
  \item $250$ см²
  \item $8$ см
  \item $288$ см³
  \item $48$ см³
  \item $60$ см²
  \item $18$ м³
  \item $180$ м³
  \item $224$ дм³
\end{enumerate}

\end{document}